%IDIOMAS (PT-BR) E CODIFICAÇÃO DE FONTES
\usepackage[T1]{fontenc}
\usepackage[utf8]{inputenc}
\usepackage[brazilian]{babel}
%OPÇÃO DE FONTES 1 (Padrão)
\usepackage{DejaVuSans}
\renewcommand*\familydefault{\sfdefault}
%OPÇAO DE FONTE 2
\usepackage{lmodern}

%DEFINIÇÃO DO VERDE OFICIAL DOS IFs
\usepackage[svgnames]{xcolor}
\definecolor{verdeOficial}{HTML}{359830}

%CLÁSSICO LOREM IPSUM
\usepackage{lipsum}

%LEGENDAS DAS IMAGENS SEM O AMBIENTE FIGURE
\usepackage{caption}

%IMAGENS
\usepackage{graphicx}

%SEM LINHA DIVISÓRIA PARA \footnote.
\renewcommand{\footnoterule}{}

%URLs
\usepackage{xurl}


%TAMANHO DO PAPEL CONFORME ABNT 15437
%Tam. mín. das margens é 2.5. Padrão para o modelo: como abaixo.
\usepackage[paperwidth=90cm,%
			paperheight=120cm,%
			top=2cm,%
			left=3.5cm,%
			right=3.5cm,%
			bottom=3.5cm,%
			]{geometry}

%TEXTO EM DUAS COLUNAS
\usepackage{multicol}
\setlength{\columnsep}{100pt}         % Distância entre as colunas
\setlength{\columnseprule}{3pt}       % Espessura da linha divisória
\renewcommand{\columnseprulecolor}{\color{verdeOficial}}	% Cor da linhas divisória

%REFERÊNCIAS E CITAÇÕES
\RequirePackage[alf,
				abnt-emphasize=bf,%
				abnt-etal-list=3,%
				abnt-etal-text=emph,%
				abnt-missing-year=sd,%
				abnt-repeated-author-omit=yes,%
				abnt-repeated-title-omit=yes,%
				abnt-thesis-year=final,%
				abnt-doi=link,%
				]{abntex2cite}
%ALTERAR O TAMANHO DA FONTE APÓS BIBLIOG. COMPILADA
\let\oldthebibliography\thebibliography
\let\endoldthebibliography\endthebibliography
\renewenvironment{thebibliography}[1]
{\normalsize\oldthebibliography{#1}}
{\endoldthebibliography}

%CAIXAS DE TEXTO PADRÃO (SEÇÕES DO PÔSTER)
\usepackage{tcolorbox}
\tcbuselibrary{many}
\tcbset{colback=green!5!white,%
		colframe=green!60!black,%
		fonttitle=\bfseries\large,%
		toprule=5mm,%
		rightrule=3mm,%
}

%CAIXA DE AGRADECIENTOS
\usepackage{awesomebox}
\newcommand{\agradecimentos}[1]{%
	\awesomebox[verdeOficial][][\textbf{\color{verdeOficial}{Agradecimentos}}]%
	{2pt}{\faThumbsUp}{verdeOficial}{#1}%
}

%DEFINIÇÃO DO RODAPÉ (LINHA)
%(RE)Fazer o cabeçalho usando apenas esse pacote, sem o minipage.
\usepackage{fancyhdr}
\renewcommand{\headrulewidth}{0.0pt}
\renewcommand{\footrulewidth}{1pt}
\renewcommand{\footrule}{%
	{\color{verdeOficial}\hrule width\headwidth height\footrulewidth}%
	\vskip\footruleskip%
}
\fancyfoot{}
